\documentclass[../DoAn.tex]{subfiles}
\begin{document}

\begin{center}
    \Large{\textbf{TÓM TẮT NỘI DUNG ĐỒ ÁN}}\\
\end{center}
\vspace{1cm}
\hspace{\parindent}Hiện nay, trong khi ngành công nghiệp game trên thế giới đang phát triển rực rỡ thì nền game Việt Nam vẫn chưa thực sự bứt phá. Một trong những rào cản lớn nhất là chi phí để sản xuất một tựa game vô cùng đắt đỏ. Thêm vào đó, chi phí bản quyền cho các công cụ phần mềm chuyên nghiệp (như hệ điều hành Windows, Photoshop, phần mềm 3D, v.v.) đang tăng cao, kết hợp cùng sự siết chặt quản lý bản quyền từ nhà nước khiến các nhà phát triển độc lập (indie) hay studio nhỏ càng gặp khó khăn trong việc vận hành và sản xuất. Để đối mặt với giới hạn kinh phí này, các hướng tiếp cận thường thấy bao gồm: làm game 2D truyền thống, phát triển game với đồ họa low-poly tối giản, hoặc phải phụ thuộc vào dòng vốn từ các nền tảng gọi vốn cộng đồng (crowdfunding) mang nhiều rủi ro.

Thay vì đi theo những phương thức đó, Đồ án này tiếp cận theo một hướng mới: tận dụng tối đa hệ sinh thái nguyên liệu (asset) chất lượng cao, giá rẻ hoặc được miễn phí hàng tuần trên nền tảng Unreal Engine 5 (UE5) để phát triển tựa game hành động 3D mang tên "Paradox Caliber". Việc lựa chọn UE5 không chỉ vì sức mạnh đồ họa mà còn xuất phát từ chính sách thân thiện với nhà phát triển độc lập: Epic Games cho phép sử dụng công cụ hoàn toàn miễn phí và không thu phí bản quyền cho tới khi sản phẩm đạt một mốc doanh thu hàng triệu đô la. Lý do cốt lõi của hướng tiếp cận này là giải quyết triệt để bài toán kinh phí khởi đầu, cho phép các đội ngũ indie an tâm sáng tạo và sở hữu được chất lượng hình ảnh đẳng cấp mà không bị cạn kiệt tài nguyên.

Tổng quan giải pháp của Đồ án không chỉ xây dựng một tựa game hoạt động trơn tru dựa trên các asset này, mà còn kiến trúc một hệ thống nhiệm vụ mở rộng. Khi kết hợp thiết kế nhiệm vụ quy chuẩn với các asset môi trường mới, nhà phát triển có thể dễ dàng thiết lập thêm một màn chơi (level) mới trong hệ thống chỉ trong thời gian rất ngắn. Từ những chương đầu về khảo sát, thiết kế và công nghệ, Đồ án sẽ làm rõ quy trình này và minh chứng cụ thể trong quá trình thực nghiệm, đóng góp nhằm giúp các nhà phát triển độc lập có thêm động lực xây dựng dự án game đồ họa hiện đại.

\begin{flushright}
Sinh viên thực hiện\\
\begin{tabular}{@{}c@{}}
\textit{(Ký và ghi rõ họ tên)}
\end{tabular}
\end{flushright}

\end{document}